\section{Finding exact certificates from partial signs}\label{sec:overview}

Consider first the feasibility problem \(Ax\le b\). The algorithm searches
for the support of a certificate. Once that support is known, the
certificate itself can be recovered by linear algebra. The difficulty is
that a coordinate essential to the certificate may be arbitrarily small
compared with the others.

\subsection{Which side contains a witness?}

Collect all possible slack vectors, together with a homogenizing
coordinate, into the subspace
\[
 C=\{(\alpha b-Ax,\alpha):x\in\RR^n,\ \alpha\in\RR\}.
\]
If \(C\) contains a nonnegative vector with positive last coordinate,
division by that coordinate gives a feasible point. Its orthogonal
complement consists of the vectors \((y,-b^\top y)\) with \(A^\top y=0\).
A nonnegative vector there with positive last coordinate is precisely a
Farkas certificate of infeasibility.

More generally, let \(C\subseteq\RR^N\) be any subspace and let
\([N]=\{1,\ldots,N\}\). Each coordinate \(i\) has a unique \emph{side}.
Its side is \(+\) if some nonnegative vector in \(C\) has positive
\(i\)th entry, and \(-\) if such a vector lies in \(C^\perp\). Farkas'
lemma gives existence, and orthogonality rules out having both
witnesses~\cite{Fark}. On the side of \(i\) one can choose a witness
that is a \emph{circuit}, a nonzero vector whose support is
inclusion-minimal among the supports of nonzero vectors in its
subspace~\cite{Rock}.

Normalize a witness to have largest entry one. An entry of size
\(2^{-T}\) can be difficult to detect even when it is indispensable.
Multiplying coordinate \(i\) of \(C\) by \(2^{-p_i}\) can bring that
entry into view. Multiplying the same coordinate of \(C^\perp\) by
\(2^{p_i}\) preserves orthogonality. We call \(p_i\) its \emph{price}.
Figure~\ref{fig:scale} illustrates why the information at a single
price vector may be incomplete.

\begin{figure}[htbp]
\centering
\begin{tikzpicture}[x=2.65cm,y=2.65cm,font=\small,
  primal/.style={lagoon,line width=1.05pt},
  dual/.style={ember,line width=.85pt,dash pattern=on 4pt off 2pt},
  axis/.style={-{Stealth[length=3pt]},line width=.45pt,black!65},
  marker/.style={fill=lagoon,draw=white,line width=.55pt}]
  \newcommand{\priceaxes}{%
    \fill[black!3] (0,0) rectangle (1,1);
    \draw[black!27,line width=.4pt] (0,1)--(1,1)--(1,0);
    \draw[axis] (-.19,0)--(1.22,0);
    \draw[axis] (0,-.2)--(0,1.24);
    \node[anchor=east,inner sep=2pt] at (-.015,1){\(1\)};
    \node[anchor=north,inner sep=2pt] at (1,-.105){\(1\)};
  }
  \begin{scope}
    \priceaxes
    \node at (.55,1.55){Original prices \(p=(0,0)\)};
    \draw[dual] (-.17,.010625)--(1.19,-.074375);
    \draw[primal] (-.01125,-.18)--(.075625,1.21);
    \path[marker] (.0625,1) circle[radius=2.3pt];
    \node[anchor=west,inner sep=2pt] at (.12,.87){\((1/16,1)\)};
    \node at (.53,-.40){\(z_1\) free,\quad \(z_2=+\)};
  \end{scope}
  \begin{scope}[xshift=7.2cm]
    \priceaxes
    \node at (.55,1.55){New prices \(p=(-4,0)\)};
    \draw[dual] (-.17,.17)--(.19,-.19);
    \draw[primal] (-.16,-.16)--(1.17,1.17);
    \path[marker] (1,1) circle[radius=2.3pt];
    \node[anchor=east,inner sep=2pt] at (.93,1.09){\((1,1)\)};
    \node at (.53,-.40){\(z_1=+,\quad z_2=+\)};
  \end{scope}
  \draw[-{Stealth[length=4pt]},black!65] (1.37,.55)--(2.15,.55);
\end{tikzpicture}
\caption{Making a small witness coordinate visible. The solid line is
the scaled primal space for \(C=\operatorname{span}\{(1/16,1)\}\), and
the dashed line is its orthogonal complement. The dot is the
nonnegative primal vector of sup norm one. Multiplying its first
coordinate by \(16\) and the corresponding dual coordinate by
\(1/16\) preserves orthogonality. At the new prices both coordinates
reach the visibility threshold \(1/2\), so both labels are forced.
Coordinates \(1\) and \(2\) are horizontal and vertical, respectively.}
\label{fig:scale}
\end{figure}


At a given \(p\), imagine a test that labels \(i\) with the correct
side whenever some scaled, normalized nonnegative vector has \(i\)th
entry at least \(1/2\). The other labels may be arbitrary. How could
these answers decide \emph{every} side? Fix a nonnegative circuit \(u\)
on side \(h\in\{-1,+1\}\). For each \(i\) in its support, the
\emph{slot} \(s=(i,h)\) records the coordinate and its side. Assign it
an integer \emph{demand}
\[
 a_s=\lfloor\log_2 u_i\rfloor+b_u,
\]
with a common shift \(b_u\) chosen so that all demands are positive.
At any slot maximizing \(a_s-hp_i\), the scaled entry is within a
factor of two of the largest scaled entry of \(u\). The label is
therefore correct at every such maximizer, even though the test
reveals neither the circuit nor the maximizing slots.

\begin{theorem}[Extremal prices]\label{thm:prices}
Let \(N,R\ge1\) be integers and let \(\Slots=[N]\times\{-1,+1\}\).
There is a deterministic algorithm that queries price vectors
\(p\in\QQ^N\), receives a label vector \(z\in\{-1,+1\}^N\) for each
query, and returns \(x\in\QQ^N\) with the following guarantee.

A hypothesis consists of a nonempty set \(S\subseteq\Slots\) and
integer demands \(a_s\in[-R,R]\) for \(s\in S\). Its surplus at
\(s=(i,h)\) is \(d_s(p)=a_s-hp_i\). Call the hypothesis
\emph{consistent} if at every query and answer \((p,z)\),
\[
   s=(i,h)\in S,\qquad
   d_s(p)=\max_{u\in S}d_u(p)\ge0
   \quad\Longrightarrow\quad z_i=h.
\]
For each consistent hypothesis the output meets all its demands,
that is, \(h x_i\ge a_s\) for every \(s=(i,h)\in S\).
The bit time, including all query and answer bits, is polynomial in
\(N\) and \(\log(R+1)\).
\end{theorem}

For the ideal test, every nonnegative circuit defines a consistent
hypothesis as soon as \(R\) bounds its demands. Since each coordinate
has a witness circuit and its demand is positive, the sign of \(x_i\)
gives its side. The prices returned by the theorem need only encode
these signs. The separate queries have already supplied the evidence.

\subsection{How to answer a price query}

Two numerical obstacles remain. A circuit can have entry ratios of
arbitrary magnitude, so the bound \(R\) from the ideal test can be
arbitrarily large at fixed dimension. Even when a rational price
\(p_i\) has few bits, the scaling factor \(2^{p_i}\) may require
prohibitively many bits to represent. The calibration procedure in
Theorem~\ref{thm:cal} gives the search a short description of the
geometry it needs. Figure~\ref{fig:architecture} shows how a query is
answered.

\begin{figure}[htbp]
\centering
\begin{tikzpicture}[x=1cm,y=1cm,font=\small,
   operation/.style={draw=black!60,fill=white,line width=.45pt,
     rounded corners=2pt,align=center,inner sep=5pt,minimum height=1.1cm},
   data/.style={-{Stealth[length=4pt]},lagoon,line width=.8pt,
     rounded corners=3pt},
   again/.style={-{Stealth[length=4pt]},ember,dashed,line width=.8pt,
     rounded corners=3pt},
   msg/.style={font=\footnotesize,inner sep=2pt,fill=white}]
  \node[operation,minimum width=2.7cm] (learn) at (1.45,0)
    {Price search\\[-1pt] \footnotesize Section~\ref{sec:prices}};
  \node[operation,minimum width=2.6cm] (cal) at (6.0,0)
    {Calibration\\[-1pt] \footnotesize Section~\ref{sec:calibration}};
  \node[operation,minimum width=3.55cm,fill=black!5] (comp) at (12.1,0)
    {Compression of\\minor magnitudes\\[-1pt] \footnotesize Section~\ref{sec:compression}};
  \node[operation,minimum width=2.9cm] (labels) at (6.0,-2)
    {Label tests by LP\\[-1pt] \footnotesize Section~\ref{sec:lp}};
  \draw[data] (learn.east)--node[msg,midway,above]{prices \(p\)} (cal.west);
  \draw[data] ([yshift=6pt]cal.east)--
     node[msg,midway,above]{requested minor} ([yshift=6pt]comp.west);
  \draw[data] ([yshift=-6pt]comp.west)--
     node[msg,midway,below]{minor sign, \(f_u\)}
     ([yshift=-6pt]cal.east);
  \draw[data] (cal.south)--node[msg,midway,right]{basis estimate \(B\)} (labels.north);
  \draw[data] (labels.west)--node[msg,above,midway]{labels \(z\)}
      (1.45,-2)--(learn.south);
  \draw[again] (comp.north)--(12.1,1.65)--
      node[msg,midway,above]{on finding a new independent minor, enlarge the set and restart}
      (1.45,1.65)--(learn.north);
\end{tikzpicture}
\caption{Answering queries during a side decision. The shaded step uses the original scalar entries. For a nonzero minor
it supplies the sign and a short substitute \(f_u\) for the logarithm
of its magnitude. The other steps use rational data of polynomial bit
length. Each restart clears the price search and retains the enlarged
set of minor indices. Once the search finishes, the signs of its
returned prices give the sides.}
\label{fig:architecture}
\end{figure}


Apply calibration to a row basis \(B_0\) of \(W=C\times C^\perp\),
with one column for each slot. Between successive restarts, the
queries share a single assignment of positive integer demands to
every circuit of \(W\). The bit length of these demands is bounded by
a polynomial in the dimension and the chosen precision. At each
queried \(p\), calibration returns a rational matrix \(B\). For
the analysis there is a reciprocal diagonal scaling \(D\) with two
properties. First, \(B\) closely approximates a row basis of \(WD\)
in which each vector of sup norm one has bounded coefficients.
Second, every slot of maximum surplus for a circuit has scaled
entry of magnitude at least one third of its largest scaled entry.
One \(D\) works for all circuits at this query. Its entries may be
very large, but only \(B\) has to be computed.

To get labels from \(B\), solve an ordinary bit-polynomial LP for each
coordinate. Give \(i\) label \(+\) if a bounded row combination of
\(B\) is nearly nonnegative, stays in a bounded box and has
\((i,+)\) entry at least \(1/4\). Otherwise give it label \(-\).
The gap between \(1/3\) and \(1/4\) absorbs the approximation error.
Opposing slots cannot both pass. In the true scaled space every
vector \(v\) has \(\sum_i v_{i+}v_{i-}=0\) by orthogonality. If both
tests passed, the midpoint of two passing combinations would have
large entries in both slots and only tiny negative entries elsewhere,
making this sum bounded away from zero. Section~\ref{sec:lp} makes the
error bounds explicit. A circuit of \(C\) or \(C^\perp\) embeds in
\(W\) by zeros on the other side, so all witness circuits remain
consistent through the final uninterrupted run of the price search.

\paragraph{Why minors suffice.}
Elimination and the description of circuits use the same data.
Fix all but one column in a square minor of \(B_0\), then vary the
last column. The resulting determinants give a vector in the row
space. Taking the fixed columns to be preceding pivots gives the
next elimination row. Taking independent columns outside a circuit's
support gives that circuit up to scale. Consequently the magnitudes
needed to choose pivots and assign demands are governed by the same
numbers, the logarithms \(\ell_u\) of the absolute nonzero minors.

Calibration uses substitutes \(f_u\) of short bit length. They preserve
the signs, including zeros, of integer affine forms with bounded
coefficients in a slight perturbation of the logarithms. To see why
this suffices, choose a pivot using \(f_u\) and the queried prices.
Record, on a fine grid, the scaled log differences from this pivot
for the other entries. For tiny relative magnitudes only an upper
bound is needed. Requiring prices to reproduce these intervals gives
linear inequalities. Eliminating the prices leaves affine
inequalities in the log data with bounded integer coefficients.
They hold for \(f\), hence also for the perturbed logs. A common
price vector for the latter therefore exists. In the true geometry
it realizes the same normalized cofactor entries to the accuracy
needed for \(B\) and for the visibility bound
(Section~\ref{sec:calibration}).

\paragraph{Why few scales suffice.}
There may be exponentially many nonzero minors. Section~\ref{sec:compression}
extends Frank--Tardos rounding~\cite{FT} to requests for one minor at
a time. The reason this works is a counting argument. Suppose the
subset sums of \(k\) minor logarithms all differ by more than two.
The corresponding \(2^k\) products of minors then have magnitudes that
more than double in sorted order, so their \(2^{2^k}\) subset sums
are distinct. Yet each sum is the value of an integer polynomial in
the matrix entries, of bounded degree and coefficient size.
Counting these polynomials bounds \(k\) by a polynomial in the
dimensions.

The procedure keeps a set of minor indices with this separation
property for accurate proxies of the logs. Lattice reduction~\cite{LLL}
provides a sequence of scales for the set. On a request, the procedure
tries to express the new proxy in those scales using bounded integer
digits and a small remainder. Success gives a short substitute, and
all answers since the last restart extend to an assignment for every
nonzero minor. Failure proves the index can be added to the separated
set. The search then starts again with a larger set, which the
counting bound permits only polynomially many times. Comparisons
for extracting digits use ordinary arithmetic on the inputs.
Small logarithms can be approximated directly. When a comparison
involves an exceptionally large logarithm it can instead be decided
by comparing products of powers of minors. In that case the input
itself has enough bits, on rational data, to bound the bit length of
the powers.

Once the search returns all side signs, actual witnesses can be
recovered by testing which coordinates can be set to zero while
retaining a witness at the desired coordinate. At minimal support
the restricted space is a line. Solving on that line recovers the
witness in the original scale. Applying feasibility to the recession
and primal--dual optimality systems then gives the alternatives in
Theorem~\ref{thm:lp}.

\subsection{Learning all demands from extremal labels}

It remains to explain how Theorem~\ref{thm:prices} turns labels into
prices. First suppose we can afford one expert per hypothesis. Give
each expert an initial unit of wealth and let it bet on a slot of
greatest surplus, with the option of abstaining at utility zero.
Choose prices to minimize the wealth-weighted sum of maximum
utilities. At this minimum, bets on opposite signs of each
coordinate balance if ties are split suitably. Multiplying wealth
according to agreement with the labels can therefore keep total
wealth from increasing, even with ties resolved against each expert.
An expert for a consistent hypothesis never loses wealth. Whenever
any of its demands is unmet, every maximizing slot has a true label
and the expert's wealth grows by a fixed factor. This can happen only
logarithmically many times in the number of experts. A majority of
queries can thus meet each consistent hypothesis's demands, and
taking a coordinatewise median preserves each such bound. The choice
of prices is related to the balanced points of Chen, Li and
Yannakakis~\cite{CLY}.

The full pool is too large to process. An expert who guesses the
demands on just four slots, however, may miss the single largest
surplus and receive no useful label. The recursive search in
Section~\ref{sec:prices} amplifies the feedback, building on
Soleymani and Farina~\cite{SF} and on earlier algorithms for
games~\cite{OAIDe,OAIRa,OAITu}. From a proposal \(y\), make a child
call that returns prices \(x\). Use the sign of \(x_i-y_i\) as
feedback at \(y\), and query the label at \(x\) to break ties.
Suppose the child has brought every surplus down to a threshold
\(t'\ge0\). For any slot \(s=(i,h)\) with \(d_s(y)\ge t'\),
\[
 d_s(x)\le t'\le d_s(y)\qquad\Longrightarrow\qquad h x_i\ge h y_i.
\]
Even equality gives correct feedback, since that slot then maximizes
surplus at \(x\). A single child can therefore give correct feedback
on every slot above a common threshold.

The analysis assigns a weight distribution on slots and a threshold
to each node for each consistent hypothesis. If few slots by weight
violate the current threshold, the child retains it. When many
violate it, the child takes a higher threshold at a surplus quantile.
Success then gives correct feedback on an upper fraction of all
slots, which four sampled slots have a constant chance to hit.
Experts guess the threshold from the price history as well as the
four demands. A median combines the children, and increasing the
weight of persistently violated slots limits the necessary depth to
\(O(1+\log N)\).

Polynomial running time also requires controlling the number of
children. An expert invests a fraction of its wealth in a child
while using the rest to bet locally. A failed child must return enough
wealth to some experts to compensate for losses on its feedback.
This accounting allows a constant number of children below the root
and rules out failure at the root. Finally, experts guessing the same
four slots and the same threshold pivot can be grouped by their
at most five demand values. The regions on which those experts act
alike have fixed dimension, so their collective bets can be computed
in polynomial bit time.

\paragraph{Other uses of the price theorem.}
For monotone linear complementarity, the signs identify which member
of each complementary pair must vanish for a solution on the
selected support. A convex quadratic test supplies labels, after
which linear feasibility gives the solution
(Section~\ref{sec:quadratic}). For a contraction in the max norm, a
coordinate of greatest fixed-point error moves toward the fixed
point under the map. This gives extremal labels from function values
(Section~\ref{sec:contractions}). Discounted value equations for
turn-based games furnish piecewise-affine examples. Section~\ref{sec:games}
uses these equations together with LP policy evaluation to recover
exact values and strategies.
